\section*{الفصل \ref{ch:g10:refraction} --- انكسار الضوء}

\begin{solution}{exo:g10:refraction:1}
تُقاس الزوايا من \omterm{def:g10:refraction:angles}{الناظم}: \omterm{def:g10:refraction:angles}{فزاوية الورود}
$90^\circ - 25^\circ = 65^\circ$، ومن ثم زاوية الانعكاس هي أيضًا
$65^\circ$. وتصنع الحزمتان الواردة والمنعكسة زاوية
$2 \times 65^\circ = 130^\circ$ فيما بينهما.
\end{solution}

\begin{solution}{exo:g10:refraction:2}
\emph{1.} $v = c/n$: ففي الماء
$v = \qty{3.00e8}{m/s}/1.33 = \qty{2.26e8}{m/s}$؛ وفي الماس
$v = \qty{3.00e8}{m/s}/2.42 = \qty{1.24e8}{m/s}$.

\emph{2.} $n = \dfrac{c}{v} = \dfrac{\qty{3.00e8}{m/s}}
{\qty{1.97e8}{m/s}} = 1.52$: أي الزجاج.
\end{solution}

\begin{solution}{exo:g10:refraction:3}
$\sin i_2 = \dfrac{\sin 40^\circ}{1.50} = \dfrac{0.643}{1.50} = 0.429$،
ومنه $i_2 = 25.4^\circ$. وبعكوسية مسارات الضوء، فإن شعاعًا داخل
الزجاج يصيب السطح بزاوية $25.4^\circ$ يخرج إلى الهواء بزاوية
$40^\circ$ بالضبط.
\end{solution}

\begin{solution}{exo:g10:refraction:4}
$\sin i_2 = \dfrac{\sin 55^\circ}{1.33} = \dfrac{0.819}{1.33} = 0.616$،
ومنه $i_2 = 38.0^\circ$. فينحني الشعاع \emph{نحو} \omterm{def:g10:refraction:angles}{الناظم}، لأن
الماء هو الوسط الأبطأ ($n_2 > n_1$).
\end{solution}

\begin{solution}{exo:g10:refraction:5}
$n = \dfrac{\sin 45^\circ}{\sin 32^\circ} = \dfrac{0.707}{0.530}
= 1.33$: فالسائل ماء.
\end{solution}

\begin{solution}{exo:g10:refraction:6}
\emph{1.} النسب الخمس هي
\[
\frac{0.342}{0.233} = 1.47,\quad
\frac{0.500}{0.334} = 1.50,\quad
\frac{0.643}{0.431} = 1.49,\quad
\frac{0.766}{0.515} = 1.49,\quad
\frac{0.866}{0.581} = 1.49 :
\]
ثابتة في حدود خطأ القياس --- أي \omterm{thm:g10:refraction:snell}{قانون سنيل}.

\emph{2.} $n \approx 1.49$: أي البلكسيغلاس.

\emph{3.} يتنبأ \omterm{thm:g10:refraction:snell}{قانون سنيل} بأن $\sin i_2 = \sin i_1 / n$: أي أن المعطيات، جيبًا
بدلالة جيب، يجب أن تصطفّ على خط مستقيم يمرّ بالمبدأ، وهو ما
تتحقق منه المسطرة بنظرة واحدة. أما رسم الزاوية بدلالة الزاوية فيضع المعطيات على
منحنٍ؛ وعند الزوايا الصغيرة يكون المنحني قريبًا من المستقيم قربًا
خادعًا، وهو بالضبط ما خدع بطليموس.
\end{solution}

\begin{solution}{exo:g10:refraction:7}
$\sin i_c = n_2/n_1$ مع $n_2 = 1.00$:
الماء $\sin i_c = 1/1.33 = 0.752$، $i_c = 48.8^\circ$؛
الزجاج $\sin i_c = 1/1.50 = 0.667$، $i_c = 41.8^\circ$؛
الماس $\sin i_c = 1/2.42 = 0.413$، $i_c = 24.4^\circ$.

ومن الهواء إلى الماء، $\sin i_2 = \sin i_1/1.33 \leq 1/1.33 < 1$ مهما تكن
\omterm{def:g10:refraction:angles}{زاوية الورود}: فثمة دائمًا زاوية \omterm{def:g10:refraction:refraction}{انكسار}، ولا يحدث شيء
خاص أبدًا --- إذ يقتضي \omterm{def:g10:refraction:critical}{الانعكاس الكلي الداخلي} $n_1 > n_2$.
\end{solution}

\begin{solution}{exo:g10:refraction:8}
\emph{1.} $d' = \dfrac{d}{n} = \dfrac{\qty{2.0}{m}}{1.33}
= \qty{1.5}{m}$.

\emph{2.} تبدو السمكة على $\qty{40}{cm}/1.33 = \qty{30}{cm}$: فالصورة
\emph{فوق} السمكة، ومن ثم على مالك الحزين أن ينقضّ \emph{تحت}
الصورة التي يراها. أما مالك الحزين، وهو غير مثقل \omterm{thm:g10:refraction:snell}{بقانون سنيل}، فيتعلم ذلك
بالتجربة والخطأ.
\end{solution}

\begin{solution}{exo:g10:refraction:9}
\emph{1.} $\sin r = \dfrac{\sin 50^\circ}{1.50} = \dfrac{0.766}{1.50}
= 0.511$، ومنه $r = 30.7^\circ$.

\emph{2.} من الزجاج إلى الهواء عند $30.7^\circ$:
$\sin i_2 = 1.50 \times 0.511 = 0.766$، ومنه $i_2 = 50^\circ$ --- أي أن
زاوية الخروج تساوي زاوية الدخول.

\emph{3.} يُلغي الانكساران أحدهما الآخر: فيغادر الشعاع
\emph{موازيًا} لاتجاهه الأصلي، مزاحًا جانبيًّا لا غير.
ولهذا لا تشوّه النافذة المشهد، إنما تزيحه إزاحة
(غير محسوسة).
\end{solution}

\begin{solution}{exo:g10:refraction:10}
$\sin 55^\circ = 0.819$. الأحمر:
$\sin i_2 = 0.819/1.61 = 0.509$، $i_2 = 30.6^\circ$. الأزرق:
$\sin i_2 = 0.819/1.66 = 0.493$، $i_2 = 29.6^\circ$. فيفصل بين اللونين
$1.0^\circ$؛ والأزرق، ذو القرينة الأكبر، ينحني أكثر
وينتهي أقرب إلى \omterm{def:g10:refraction:angles}{الناظم}.
\end{solution}

\begin{solution}{exo:g10:refraction:11}
\emph{1.} $\sin i_c = \dfrac{1.46}{1.48} = 0.986$، ومنه
$i_c = 80.6^\circ$.

\emph{2.} شعاع بزاوية $i_c$ عن \omterm{def:g10:refraction:angles}{الناظم} يصنع $90^\circ - i_c$ مع
المحور، فيكون مساره على طول محور طوله $L$ هو $\dfrac{L}{\sin i_c}$:
\[
\qty{1000}{m} \times \frac{1.48}{1.46} = \qty{1013.7}{m},
\]
أي نحو \qty{14}{m} زيادةً على الشعاع المحوري.
\end{solution}

\begin{solution}{exo:g10:refraction:12}
\emph{1.} الضوء الملامس الآتي من الأفق يصل عند $i_1$ قريبة من
$90^\circ$؛ فينكسر إلى $\sin i_2 = \sin 90^\circ / 1.33 = 0.752$،
أي\ $i_2 = 48.8^\circ$ --- أي \omterm{def:g10:refraction:critical}{الزاوية الحدية}، كما تقتضي
العكوسية. ولذلك تصل أشعة السماء كلها إلى العين داخل مخروط
نصف زاويته $48.8^\circ$.

\emph{2.} حافة القرص هي حيث يلاقي ذلك المخروط السطح:
\[
r = d \tan i_c = \qty{2.0}{m} \times \tan 48.8^\circ
= \qty{2.0}{m} \times 1.14 = \qty{2.3}{m}.
\]

\emph{3.} خارج القرص، يصيب الضوء الآتي من تحت السطحَ بزاوية أبعد
من \omterm{def:g10:refraction:critical}{الزاوية الحدية} فينعكس كليًّا: فيرى الغطّاس مرآة
فضية تعرض قاع المسبح.
\end{solution}

\begin{solution}{exo:g10:refraction:13}
\emph{1.} الأحمر: $\sin i_2 = 1.51 \times \sin 30^\circ = 0.755$، ومنه
$i_2 = 49.0^\circ$. الأزرق: $\sin i_2 = 1.53 \times 0.500 = 0.765$، ومنه
$i_2 = 49.9^\circ$. فتنتشر الحزمة البيضاء على
$49.9^\circ - 49.0^\circ = 0.9^\circ$ --- أي \omterm{def:g10:light-spectra:white}{طيف} موشور من \omterm{def:g10:refraction:refraction}{انكسار}
واحد.

\emph{2.} عند $45^\circ$: $\sin i_2 = 1.51 \times 0.707 = 1.07 > 1$
(و$1.53 \times 0.707 = 1.08$ للأزرق). فلا وجود لزاوية \omterm{def:g10:refraction:refraction}{انكسار}
لأي من اللونين: إذ \omterm{def:g10:refraction:critical}{الزاوية الحدية}
$\arcsin(1/1.51) = 41.5^\circ < 45^\circ$، فيعكس الوجه الطويل الحزمة كلها
انعكاسًا كليًّا.
\end{solution}

\begin{solution}{exo:g10:refraction:14}
\emph{1.} \omterm{def:g10:refraction:critical}{الزاوية الحدية} بين الزجاج والهواء: $\sin i_c = 1/1.50$، ومنه
$i_c = 41.8^\circ < 45^\circ$: فتنعكس الحزمة انعكاسًا كليًّا --- أي $100\%$
من الضوء، دون طلاء يبهت. أما المرآة المعدنية فتفقد بضعة
بالمئة عند كل ارتداد، والمِرقاب يحتاج ارتدادين.

\emph{2.} بين الزجاج والماء: $\sin i_c = \dfrac{1.33}{1.50} = 0.887$، ومنه
$i_c = 62.5^\circ$. والآن $45^\circ < i_c$: أي أن الورود دون
الحدية، فينكسر معظم الضوء خارجًا إلى الماء، وتخفق
المرآة --- ومن ثم تخفق الصورة.

\emph{3.} $\sin i_2 = \dfrac{1.50}{1.33} \times \sin 45^\circ
= 1.128 \times 0.707 = 0.798$، ومنه $i_2 = 52.9^\circ$ إلى الماء.
\end{solution}

\begin{solution}{exo:g10:refraction:15}
\emph{1.} $v = \dfrac{\qty{3.00e8}{m/s}}{1.47} = \qty{2.04e8}{m/s}$، ومنه
\[
t = \frac{\qty{6.2e6}{m}}{\qty{2.04e8}{m/s}} = \qty{30.4}{ms} \approx
\qty{30}{ms}.
\]

\emph{2.} الراديو: $t = \qty{6.2e6}{m}/\qty{3.00e8}{m/s} =
\qty{20.7}{ms}$. فالزجاج يكلّف $30.4 - 20.7 = \qty{9.7}{ms}$ في كل
اتجاه، أي نحو \qty{19}{ms} للذهاب والإياب.

\emph{3.} خفض القرينة --- فالألياف ذات \omterm{def:g10:refraction:fiber}{القلب} المجوّف توجّه الضوء عبر
الهواء ($n \approx 1.00$)، وتفعل الوصلات المكروية عبر \omterm{def:g10:refraction:fiber}{الغلاف} الجوي
الشيء نفسه --- أو تقصير المسار: أي كابل يُمدّ أقرب إلى
مسار الدائرة العظمى بين المدينتين. وكلاهما يُباع فعلًا، بأثمان
لكل جزء من الألف من \omterm{def:g10:orders-of-magnitude:unit}{الثانية}.
\end{solution}

\begin{solution}{pb:g10:refraction:1}
\textbf{1.} $v_{\text{الماء}} = \qty{3.00e8}{m/s}/1.33 =
\qty{2.26e8}{m/s}$؛ $v_{\text{الماس}} = \qty{3.00e8}{m/s}/2.42 =
\qty{1.24e8}{m/s}$ --- فالماس يبطّئ الضوء بالعامل $n = 2.42$.

\textbf{2.} $\sin i_2 = \sin 40^\circ / 1.33 = 0.643/1.33 = 0.483$، ومنه
$i_2 = 28.9^\circ$.

\textbf{3.} $\sin i_2 = 1.33 \times \sin 28.9^\circ = 1.33 \times
0.483 = 0.643$، ومنه $i_2 = 40^\circ$: فالضوء يعيد اقتفاء مساره ---
أي عكوسية الأشعة الضوئية.

\textbf{4.} $\sin i_c = 1/1.33 = 0.752$، ومنه $i_c = 48.8^\circ$. وعند
$60^\circ > i_c$، سيحتاج \omterm{thm:g10:refraction:snell}{قانون سنيل} إلى $\sin i_2 = 1.33 \sin
60^\circ = 1.15 > 1$: أي لا \omterm{def:g10:refraction:refraction}{شعاع منكسر}؛ فيعكس السطح
الضوء كليًّا راجعًا إلى أسفل.

\textbf{5.} $v = \qty{3.00e8}{m/s}/1.0003 = \qty{2.999e8}{m/s}$: فالضوء
في الهواء أبطأ منه في الخلاء بمقدار $0.03\%$، وهو أدنى بكثير من دقة
قياسات زوايانا، ولذلك فإن التقريب $n_{\text{الهواء}} = 1.00$ غير ضارّ.

\textbf{6.} $d' = \qty{3.0}{m}/1.33 = \qty{2.3}{m}$: فالمسبح يخفي
\qty{70}{cm} من نفسه.

\textbf{7.} ينكسر شعاعان يغادران قطعة النقود عند السطح، فينحنيان
\emph{مبتعدين} عن \omterm{def:g10:refraction:angles}{الناظم}، ومن ثم يتباعدان فوق الماء تباعدًا أشدّ
انحدارًا مما تحته. فتمدّ العين الشعاعين الخارجين رجوعًا في
خطوط مستقيمة؛ فيتقاطعان \emph{فوق} قطعة النقود، وذلك التقاطع هو
حيث تضع العين الصورة.

\textbf{8.} $r = d \tan i_c = \qty{3.0}{m} \times \tan 48.8^\circ =
\qty{3.0}{m} \times 1.14 = \qty{3.4}{m}$: أي قرص قطره نحو \qty{6.9}{m}
يسع السماء كلها.

\textbf{9.} عند $30^\circ$: $\sin i_2 = 1.33 \times 0.500 = 0.665$، فيخرج
الشعاع عند $41.7^\circ$. وعند $45^\circ$: $\sin i_2 = 1.33 \times
0.707 = 0.940$، فيخرج عند $70.1^\circ$، ملامسًا تقريبًا. وعند
$50^\circ$: $1.33 \times 0.766 = 1.02 > 1$ --- أي انعكاس كلي
داخلي، فيغوص الشعاع راجعًا إلى المسبح.

\textbf{10.} الصورة فوق السمكة الحقيقية
(\cref{prop:g10:refraction:depth})، فصوّب \emph{تحتها}. ونعم:
ترى السمكة الصيادَ مزاحًا هو أيضًا، مضغوطًا نحو العمودي
داخل نافذة سنيل --- فكل يرى الآخر حيث ليس الآخر.

\textbf{11.} الماس: $\sin i_c = 1/2.42 = 0.413$، ومنه
$i_c = 24.4^\circ$، مقابل $41.1^\circ$ للزجاج: فمخروط إفلات الماس
لا يبلغ نصف عرض مخروط الزجاج.

\textbf{12.} في الماس، $30^\circ > 24.4^\circ$: أي انعكاس كلي
داخلي، فيبقى الشعاع في الحجر. وفي الزجاج،
$30^\circ < 41.1^\circ$: فيتسرّب الشعاع من الخلف. أما أوجه
القطع البريليانتي فتُمال بحيث يصيبها الضوء الداخل دائمًا تقريبًا
بزاوية أبعد من $24.4^\circ$ فيرتدّ حتى يخرج صاعدًا نحو
العين --- أما التقليد الزجاجي، بمخروط إفلاته الواسع، فيدع الضوء
يسقط من خلاله ببساطة.

\textbf{13.} $\sin 45^\circ = 0.707$. الأحمر: $0.707/2.41 = 0.293$، ومنه
$i_2 = 17.1^\circ$. الأزرق: $0.707/2.45 = 0.289$، ومنه $i_2 = 16.8^\circ$.
والفصل عند الدخول: نحو $0.3^\circ$.

\textbf{14.} $\sin 17.0^\circ = 0.292$. الأحمر: $\sin i_2 = 2.41 \times
0.292 = 0.705$، ومنه $i_2 = 44.8^\circ$. الأزرق: $\sin i_2 = 2.45 \times
0.292 = 0.716$، ومنه $i_2 = 45.8^\circ$. فعرض المروحة الآن $1.0^\circ$:
إذ إن \omterm{def:g10:refraction:refraction}{انكسار} الخروج، وهو يشتغل قرب \omterm{def:g10:refraction:critical}{الزاوية الحدية}، مدّد
فصل الدخول $0.3^\circ$ ثلاثة أضعاف. وذلك الرذاذ الممدَّد من الألوان
هو \emph{نار} الماسة.

\textbf{15.} بين الماس والماء: $\sin i_c = 1.33/2.42 = 0.550$، ومنه
$i_c = 33.3^\circ$. فالشعاع $30^\circ$، وقد كان محبوسًا حين كان الوجه مواجهًا
للهواء، يفلت الآن ($30^\circ < 33.3^\circ$) إلى غشاء الماء: فكل
قطرة أو بصمة أو غشاء دهن يرفع القرينة الخارجية، فيوسّع
مخروط الإفلات ويستنزف البريق --- ومن هنا خرقة الجواهري.

\textbf{16.} $\sin i_c = \dfrac{1.46}{1.48} = 0.986$، ومنه
$i_c = 80.6^\circ$.

\textbf{17.} يحدث \omterm{def:g10:refraction:critical}{الانعكاس الكلي الداخلي} \emph{عند السطح}،
فوجب أن يبقى السطح تامًّا: ففي خيط عارٍ، ترفع كل ذرة
غبار أو خدش أو قطرة ماء أو إصبع لامسة القرينةَ الخارجية
عند تلك النقطة فتنزف الضوء. أما \omterm{def:g10:refraction:fiber}{الغلاف} فيدفن المرآة
داخل زجاج نظيف لا يبلغه شيء --- ويحمي، فوق ذلك،
القلبَ الرفيع كالشعرة حمايةً ميكانيكية.

\textbf{18.} $v = \qty{3.00e8}{m/s}/1.48 = \qty{2.03e8}{m/s}$، ومنه
\[
t = \frac{\qty{6.0e6}{m}}{\qty{2.03e8}{m/s}} = \qty{29.6}{ms} \approx
\qty{30}{ms}.
\]

\textbf{19.} عامل المسار $1/\sin i_c = 1.48/1.46 = 1.0137$: أي في أسوأ الأحوال
$\qty{6000}{km} \times 0.0137 \approx \qty{82}{km}$ من الزجاج الزائد،
تكلّف $\qty{8.2e4}{m} / \qty{2.03e8}{m/s} \approx \qty{0.4}{ms}$ ---
فالتعرّج شبه مجاني.

\textbf{20.} البِنغ: $2 \times (29.6 + 0.4) \approx \qty{60}{ms}$.
والصبيب: $\dfrac{10^{13}}{2.4 \times 10^8} \approx 40\,000$ نسخة
من هذا الكتاب في \omterm{def:g10:orders-of-magnitude:unit}{الثانية}. وفي جملة واحدة: بحبس الضوء داخل
الزجاج \omterm{def:g10:refraction:critical}{بالانعكاس الكلي الداخلي}، يحمل \omterm{thm:g10:refraction:snell}{قانون سنيل} مكتبةً
عبر المحيط كل \omterm{def:g10:orders-of-magnitude:unit}{ثانية}، ستين جزءًا من الألف من \omterm{def:g10:orders-of-magnitude:unit}{الثانية} ذهابًا وإيابًا.
\end{solution}
